% Formatvorlage f�r studentische Arbeiten der Arbeitsgruppe Software Systems
% Engineering am Institut f�r Informatik der Universit�t Hildesheim
%   erstellt von Christopher Voges, 28.04.2016
%   �berarbeitet von Sascha El-Sharkawy, 06.08.2019
%   https://sse.uni-hildesheim.de/studium-lehre/richtlinien-fuer-ausarbeitungen/vorlagen/

% Dokumentenkopf 
\documentclass[
    11pt, % Schriftgr��e
    DIV=10,
    ngerman, % f�r Umlaute, Silbentrennung etc.
    a4paper, % Papierformat
    twoside, % weiseitiges Dokument
    titlepage, % es wird eine Titelseite verwendet
    parskip=half, % Abstand zwischen Abs�tzen (halbe Zeile)
    headings=normal, % Gr��e der �berschriften verkleinern
    listof=totoc, % Verzeichnisse im Inhaltsverzeichnis auff�hren
    bibliography=totoc, % Literaturverzeichnis im Inhaltsverzeichnis auff�hren
    index=totoc, % Index im Inhaltsverzeichnis auff�hren
    captions=tableheading, % Beschriftung von Tabellen unterhalb ausgeben
		numbers=noenddot,
    final, % Status des Dokuments (final/draft)
]{scrreprt}
\usepackage{subfiles}

% Einbinden der Packages
% Entfernt Probleme durch Verwendung von Koma-Script Klassen (scrreport)
\usepackage{scrhack}
% Erlaubt Konfiguration von Header & Footer
\usepackage[automark,headsepline,footsepline,plainheadsepline,plainfootsepline]{scrlayer-scrpage}

% Anpassung an Sprache (deutsch)
\usepackage[ngerman]{babel}

% Umlaute
\usepackage[latin1]{inputenc}
\usepackage[T1]{fontenc}
\usepackage{textcomp} % Euro-Zeichen etc.

% Schrift
\usepackage{lmodern}
\usepackage{relsize}

% Einbinden von JPG-Grafiken erm�glichen
\usepackage[dvips,final]{graphicx}

% Erm�glichen mathematischer Symbole
\usepackage{amsmath,amsfonts}

% F�r die Definition der Zeilenabst�nde, Seitenr�nder etc.
\usepackage{setspace}
\usepackage{geometry}

% URL-Unterst�tzung
\usepackage{url}

% Abk�rzungsverzeichnis 
% Alles weitere hierzu in: "Inhalt\Abkuerzungen.tex".
\usepackage[intoc]{nomencl}
\let\abbrev\nomenclature
\renewcommand{\nomname}{Abk�rzungsverzeichnis}
\setlength{\nomlabelwidth}{.15\textwidth}

% PDF-Optionen -----------------------------------------------------------------
\usepackage[
    bookmarks, % Es werden Bookmarks verwendet
    bookmarksopen=true, % Farbe von Bookmarks
    colorlinks=true, % Farbe von Verkn�pfungen
    linkcolor=black, % einfache interne Verkn�pfungen
    anchorcolor=black, % Ankertext
    citecolor=black, % Verweise auf Literaturverzeichniseintr�ge im Text
    filecolor=black, % Verkn�pfungen, die lokale Dateien �ffnen
    menucolor=black, % Acrobat-Men�punkte
    urlcolor=black, % Farbe der URLs
    plainpages=false, % zur korrekten Erstellung der Bookmarks
    pdfpagelabels, % zur korrekten Erstellung der Bookmarks
    hypertexnames=false, % zur korrekten Erstellung der Bookmarks
    linktocpage, % Seitenzahlen anstatt Text im Inhaltsverzeichnis verlinken
    pdfusetitle % Erm�glicht das Setzen der Meta-Daten des erzeugten PDFs
]{hyperref}
% \renewcommand{\theHsection}{\thepart.section.\thesection}
% \hypersetup{
%     %pdftitle={\titel \untertitel},
%     pdfauthor={\autor},
%     pdfcreator={\autor}
%     %pdfsubject={\titel \untertitel},
%    % pdfkeywords={\titel \untertitel}
% }

% Wird f�r Teile der Formatierung des Deckblatts und die Verwendung von
% Aufz�hlungen ben�tigt
\usepackage{listings}
\usepackage{xcolor} 


% fortlaufendes Durchnummerieren der Fu�noten
\usepackage{chngcntr}

% bei der Definition eigener Befehle ben�tigt
\usepackage{ifthen}

% sorgt daf�r, dass Leerzeichen hinter parameterlosen Makros nicht als Makroendezeichen interpretiert werden
\usepackage{xspace}

% F�r das Erstellen eines Glossars
\usepackage[toc, automake, nonumberlist]{glossaries}

\usepackage[figure,table,lstlisting]{totalcount}

% Einbinden der Meta-Informationen
% Meta-Informationen 
% Falls Umlaute oder ein *�* vorkommen:
\usepackage[latin1]{inputenc}

% Hier k�nnen Sie Informationen zur Arbeit, sich selbst und Ihren Betreuern
% hinterlegen.



\newcommand{\titel}{Search-Based Generation of Test Data - Ein praktischer �berblick (DRAFT)} % Name der Arbeit
% Test Data Generation Tools - the Search-Based Way
% Search-Based Test Data Generation Tools

\newcommand{\untertitel}{} % Optional mit Untertitil
% Art der Arbeit ggf. zus�tzlich der Titel der Veranstaltung
\newcommand{\art}{Seminararbeit}
\newcommand{\studiengang}{Angewandte Informatik (BSc)} % Ihr Studiengang
\newcommand{\autor}{Dennis Pidun} % Ihr Name
\newcommand{\email}{pidund@uni-hildesheim.de}% Ihre aktuelle und g�ltige E-Mail-Adresse
\newcommand{\matnr}{296267} % Ihre Matrikelnummer

% Die Angaben zu uns:
\newcommand{\institut}{Institut f�r Informatik}
\newcommand{\arbeitsgruppe}{Arbeitsgruppe Software Systems Engineering}
\newcommand{\erstgutachter}{Prof. Dr. Klaus Schmid, SSE}% Ihr(e) ErstgutachterIn
\newcommand{\zweitgutachter}{Adam Krafczyk (MSc)}% Ihr(e) ZweitgutachterIn
\newcommand{\universitaet}{Universit�t Hildesheim\ \textbullet \ Universit�tsplatz 1 \ \textbullet \ D-31134 Hildesheim}
\newcommand{\adresse}{\arbeitsgruppe \ \textbullet \ \institut \\ \universitaet}

\newcommand{\version}{Version 0.0.1}% Die Version der Arbeit

% Wird 'projektarbeit' auf 'true' gesetzt, wird keine Eigenst�ndigkeitserkl�rung
% erzeugt.
\newboolean{projektarbeit}
\setboolean{projektarbeit}{true}

% Wird 'final' auf 'true' gesetzt, werden folgende �nderungen vorgenommen:
% -Entfernen von Datum in der Kopf- und Versionsnummer in der Fu�zeile
% -Entfernen von Datum und Versionsnummer vom Deckblatt
% -Es werden Leerseiten f�r den doppelseitigen Druck eingef�gt
\newboolean{final}
\setboolean{final}{false}





% Erstellung der Verzeichnisse und Glossars aktivieren
\makeindex
\makenomenclature
\glstoctrue 


% Kopf- und Fu�zeilen, Seitenr�nder etc. anpassen
% Zeilenabstand: einfach 
\newcommand{\zeilenabstandHauptteil}{1.0}
\newcommand{\zeilenabstandAnhang}{1.0}

% Seitenr�nder
\setlength{\topskip}{\ht\strutbox} % behebt Warnung von geometry
% Initiales Papierformat, wird nach dem Deckblatt teilwiese ge�ndert 
\geometry{paper=a4paper,left=3.5cm, right=2.5cm, top=2.5cm, bottom=2.9cm, headsep=.6cm}

%% Einstellen der Schriftgr��en der �berschriften
\setkomafont{chapter}{\LARGE\bfseries\rmfamily}
\setkomafont{section}{\Large\bfseries\rmfamily}
\setkomafont{subsection}{\large\bfseries\rmfamily}
\setkomafont{subsubsection}{\large\mdseries\rmfamily}
\setkomafont{pageheadfoot}{\normalfont}

%% Einstellen der Abst�nde vor und nach den �berschriften
\RedeclareSectionCommand[
  %runin=false,
  afterindent=false,
  beforeskip=0pt,
  afterskip=.05\baselineskip]{chapter}
\RedeclareSectionCommand[
  %runin=false,
  afterindent=false,
  beforeskip=0pt,
  afterskip=.05\baselineskip]{section}
\RedeclareSectionCommand[
  %runin=false,
  afterindent=false,
  beforeskip=0pt,
  afterskip=.05\baselineskip]{subsection}
\RedeclareSectionCommand[
  %runin=false,
  afterindent=false,
  beforeskip=0pt,
  afterskip=.05\baselineskip]{subsubsection}

% Header & Footer konfigurieren
\renewcommand*{\chaptermarkformat}{} % Keine Nummerierung im Header
\clearpairofpagestyles % Voreinstellungen l�schen
% E = Even page  (Gerade Seitennummer)
% O = Odd page   (Ungerade Seitennummer)
% L = Left side	 (Linker Teil der Seite)
% C = Centered   (Mittlerer Teil der Seite)
% R = Right side (Rechter Teil der Seite)
% Auf ungeraden Seiten: Kapitel�berschrift, oben rechts
\rohead[\leftmark]{\leftmark}
% Auf geraden Seiten: Titel der Arbeit, oben links
\lehead[\titel]{\titel}
% Seitenzahlen in der Mitte eines Zweiseitigenausdrucks
\rofoot[\thepage]{\thepage}
\lefoot[\thepage]{\thepage}

\ifthenelse{\boolean{final}}{% Keine besondere Aktion in der finalen Fassung
  } {
		%Versionsnummer f�r Draft
		\lohead[Version vom \today~(Vor Abgabe entfernen)]{Version vom \today~(Vor Abgabe entfernen)}
		\rehead[Version vom \today~(Vor Abgabe entfernen)]{Version vom \today~(Vor Abgabe entfernen)}
		\lofoot[Version vom \version~(Vor Abgabe entfernen)]{Version vom \version~(Vor Abgabe entfernen)}
		\refoot[Version vom \version~(Vor Abgabe entfernen)]{Version vom \version~(Vor Abgabe entfernen)}
	}

% erzeugt ein wenig mehr Platz hinter einem Punkt
\frenchspacing 

% Schusterjungen und Hurenkinder vermeiden
\clubpenalty = 10000
\widowpenalty = 10000 
\displaywidowpenalty = 10000

% Quellcode-Ausgabe formatieren
\lstset{numbers=left, numberstyle=\tiny, numbersep=5pt, breaklines=true}
\lstset{emph={square}, emphstyle=\color{red}, emph={[2]root,base}, emphstyle={[2]\color{blue}}}

% Fu�noten fortlaufend durchnummerieren
\counterwithout{footnote}{chapter}

% Einstellungen f�r Listings
\definecolor{hellgelb}{rgb}{1,1,0.9}
\definecolor{colKeys}{rgb}{0,0,1}
\definecolor{colIdentifier}{rgb}{0,0,0}
\definecolor{colComments}{rgb}{1,0,0}
\definecolor{colString}{rgb}{0,0.5,0}
\lstset{
    float=hbp,
    columns=flexible,
    tabsize=2,
    frame=single,
    extendedchars=true,
    showspaces=false,
    showstringspaces=false,
    numbers=left,
    numberstyle=\tiny,
    breaklines=true,
    breakautoindent=true,
    xleftmargin=0.6cm,
    xrightmargin=0.1cm,
    captionpos=b
}

% Eigene Definitionen f�r Silbentrennung laden
% Trennvorschl�ge im Text werden mit \" angegeben
% untrennbare W�rter und Ausnahmen von der normalen Trennung k�nnen in dieser
% Datei mittels \hyphenation definiert werden


% Eigene LaTeX-Befehle laden
% Eigene Befehle, z.B.:
\newcommand{\fett}[1]{\textbf{#1}}

% Hier beginnt das eigentliche Dokument
\begin{document}
\title{\titel}
\author{\autor}

% Setzt wie tief Abschnitte numeriert und ins Inhaltsverzeichnis
% aufgenommen werden sollen, hier: bis inkl. SubSubSection (z.B. 1.1.1.1).
\setcounter{secnumdepth}{3}
\setcounter{tocdepth}{3}

% Deckblatt einbinden
% Erzeugt das Deckblatt
%   Bei zu langem Arbeitstitel m�ssen die vertikalen Abst�nde (\vspace)
%   angepasst werden, damit das Deckblatt weiterhin auf eine Seite passt.
\begin{titlepage}
\newgeometry{top=2cm,bottom=2cm,left=2cm,right=2cm}
\begin{figure}
    \includegraphics[scale=0.22]{Bilder/SSE_logo_2000_resized.pdf}
		\hfill
		\includegraphics[scale=0.25]{Bilder/St_Uni-Logo-9-2003-eps-converted-to.pdf}
\end{figure}
\begin{center}
    \vspace*{0cm}
    \huge{\textbf{\art~im Studiengang \studiengang}}
\end{center}
\begin{center}
    \vspace{1cm}
    \Huge{\textbf{\titel}}
    \vspace{1cm}
\end{center}
\ifthenelse{\boolean{final}}{}{
    \begin{center}
        \version ~vom \today\\
        (Vor Abgabe entfernen)
    \end{center}
}
\begin{center}
    \vspace{1cm}
    \huge{\textbf{\autor\\}}
    \vspace{1cm}
    \LARGE{\matnr\\}
    \vspace{0,5cm}
    \LARGE{\email}
\end{center}
\vspace{1cm}
\begin{center}
    \centering
    \LARGE{\textbf{Betreuer:}} \\
    \LARGE{\erstgutachter} \\
    \LARGE{\zweitgutachter} ~\\
\end{center}
\vspace{0.5cm}
\begin{flushright}
\end{flushright}
\begin{center}
    \small{\arbeitsgruppe \ \textbullet \ \institut \\ Universit�t Hildesheim
    \textbullet \ Universit�tsplatz 1 \textbullet \ D-31134 Hildesheim}
\end{center}
\end{titlepage}
\restoregeometry

% Setzen des Papierformats f�r den Rest des Dokuments
%\newgeometry{left=3.5cm, right=2.5cm, top=2.9cm, bottom=2.9cm}

% Bei finaler Fassung: Leere Seite als R�ckseite des Deckblatts 
\ifthenelse{\boolean{final}}{\cleardoublepage}{} 

% Selbstst�ndigkeitserkl�rung einbinden
\ifthenelse{\boolean{projektarbeit}}{}{\thispagestyle{empty}
\vspace*{.5cm}
\Large \textbf{Eigenst�ndigkeitserkl�rung} \normalsize
\begin{verbatim}
\end{verbatim}\setstretch{1.20}
% Bei abweichender Vorgabe bspw. durch Ihr Pr�fungsordnung, muss der
% nachfolgende Text durch den vorgegeben ersetzt werden!
\fett{Erkl�rung �ber das selbstst�ndige Verfassen von "\titel"}

Ich versichere hiermit, dass ich die vorstehende Arbeit selbstst�ndig verfasst
und keine anderen als die angegebenen Hilfsmittel benutzt habe. Die Stellen der obigen Arbeit, die anderen Werken dem Wortlaut oder dem Sinn nach entnommen wurden, habe ich in jedem Fall durch die Angabe der Quelle bzw. der Herkunft, auch der benutzten Sekund�rliteratur, als Entlehnung kenntlich gemacht. Dies gilt auch f�r Zeichnungen, Skizzen, bildliche Darstellungen sowie f�r Quellen aus dem Internet und anderen elektronischen Text- und Datensammlungen und dergleichen. Die eingereichte Arbeit ist nicht anderweitig als Pr�fungsleistung verwendet worden oder in deutscher oder einer anderen Sprache als Ver�ffentlichung erschienen. Mir ist bewusst, dass wahrheitswidrige Angaben als T�uschung behandelt werden.

Hildesheim, den~\today

\unitlength 5mm
\begin{picture}(5,4) \put(0,0) {\line(1,0){22}} \end{picture}\\
\autor
}

% Bei finaler Fassung: Leere Seite als R�ckseite der Erkl�rung
\ifthenelse{\boolean{projektarbeit}}{}{\ifthenelse{\boolean{final}}{\cleardoublepage}{}}

% Abstract einbinden
% Vor dem Hauptteil werden die Seiten in gro�en r�mischen Ziffern nummeriert.
\pagenumbering{roman}

\section*{Kurzfassung}
\section*{Abstract}



% Bei finaler Fassung: Leere Seite als R�ckseite des Abstracts
\ifthenelse{\boolean{final}}{\cleardoublepage}{}

% Verzeichnisse drucken
\tableofcontents% Inhaltsverzeichnis

% Wenn Abbildungs- und Tabellenverzeichnis auf die selbe Seite sollen: 
\iftotalfigures\listoffigures\fi
\begingroup 
\let\clearpage\relax
\vspace{1cm} 
\iftotaltables\listoftables\fi
\endgroup
 
% Wenn Abbildungs- und Tabellenverzeichnis je auf einer eigenen Seite beginnen
% sollen: 
% \listoffigures % Abbildungsverzeichnis
% \listoftables % Tabellenverzeichnis

% Erstellen des Liste der Listings
\renewcommand{\lstlistlistingname}{Quellcode-Verzeichnis}
\iftotallstlistings\lstlistoflistings\fi

% F�r korrekte �berschrift in der Kopfzeile
\clearpage\markboth{\nomname}{\nomname}

% Drucken des Abk�rzungsverzeichnises
\printnomenclature
\label{cha:Abkuerzungsverzeichnis}
 
% Arabische Seitenzahlen im Hauptteil 
\ifthenelse{\boolean{final}}{\cleardoublepage}{\clearpage}
\pagenumbering{arabic}

% Die Inhaltskapitel aus "Inhalt.tex" einbinden
\begin{spacing}{\zeilenabstandHauptteil}
% Hier k�nnen die einzelnen Kapitel inkludiert werden. 
% Die Dateien m�ssen auf .tex enden. Diese Endung muss
% beim Inkludieren aber weggelassen werden.
% Info: \include und \input unterscheiden sich im wesentlichen darin, dass bei
% \include immer eine neue Seite angefangen wird.

\chapter{Einleitung}
\section{Motivation}
\section{Ziel dieser Arbeit}

\chapter{Background Information}
\section{Symbolic Execution}
Zur Analyse des zu testenden Programms wird es mithilfe von symbolischen Werten ausgef�hrt.
Um dies zu demonstrieren, stelle man sich eine Funktion f(x) und eine Funktion h(x, y) vor. Zur Vereinfachung ist die
Gr��e der verwendeten Funktionen gering. Zu betrachten sei somit folgender Code-Ausschnitt in Listing \ref{lst:function-fh}.

\begin{minipage}{\linewidth}
    \lstset{language=c, basicstyle=\footnotesize, showstringspaces=false,tabsize=2}
    \lstinputlisting[label={lst:function-fh},caption={Demonstrationsbeispiel von DART nach Godefroid}]{listings/function-fh.code}
\end{minipage}

Hierbei werden f�r jede Verschachtelung equivalente Gleichungen geschaffen, welche statt Programmvariablen symbolische
Werte beinhalten. Somit folgt \(x \mapsto x_0\) und \(y \mapsto y_0\). Mithilfe dieser symbolischen Werten lassen sich also
symbolische Formeln ableiten, welche aus der Verfolgung der Unterprogramme entstehen \cite{godefroid2005dart}. Hierzu
werden fortlaufend alle Verschachtelungen untersucht und wo es m�glich ist nachverfolgt. Beispielsweise sieht man in
Listing \ref{lst:function-fh} in Zeile 3 die Bedingung \(x \neq y\). Aus der Analyse zur symbolischen Ausf�hrung folgt
daher \(x \neq y \mapsto x_0 \neq y_0\) und f�r Zeile 4 \(f(x) = x+10 \mapsto 2*x_0 = x_0+10\). Die symbolischen
Variablen \(x_0\) und \(y_0\) enthalten hier die Speicheradressen f�r \(x\) und \(y\). Somit lassen sich �ber die
Ausf�hrung mithilfe von symbolischen Werten Gleichungssysteme bestimmen, welche dann gel�st werden k�nnen \cite{godefroid2005dart}.

\chapter{Dynamic Approach}
Beim dynamischen Ansatz von Korel wird das zu testende Programm zur Laufzeit analysiert und Eingabedaten generiert.
Zun�chst wird der Ansatz von Korel dadurch eingeschr�nkt, dass dieser  sich auf das Node Problem konzentriert.
Dieses wird definiert durch:
\begin{quote}
    Given node Y in a program. The goal is to find a program input x on which node Y will be executed \cite{korel1990dynamic}.
\end{quote}

Dieses Problem wurde in der Vergangenheit damit gel�st, dass dieses auf das Path Problem reduziert wurde \cite{korel1990dynamic}.
Bei diesem Ansatz werden nacheinander alle m�glichen Wege zu eine Node Y gew�hlt. F�r jeden m�glichen Weg werden dann
entsprechend Eingabewerte generiert, die zu Node Y f�hren k�nnten. Dies ben�tigt jedoch das L�sen einer Reihe von
Ungleichungen, die in \textit{if} -und \textit{repeat} Statements vorkommen \cite{ince1987automatic}. H�ufig wird daher
das Programm mit symbolischen Werten ausgef�hrt. Diese Ausf�hrung ist jedoch in bestimmten F�llen sehr aufwendig \cite{ince1987automatic},
weshalb hier andere Ans�tze verfolgt werden sollten. Zudem kann es h�ufiger dazu kommen, dass einige Wege nicht m�glich sind,
und daher Ressourcen f�r die Berechnung m�glicher Eingabewerte verschwendet werden \cite{korel1990dynamic}.

Korel schl�gt daher einen anderen Ansatz vor, bei dem die Selektierung des Weges wegf�llt \cite{korel1990dynamic}.
Hierzu wird das Programm mit einer zuf�lligen Eingabe ausgef�hrt. Bei jedem neuen Branch wird eine \textit{search procedure}
angesto�en, welche entscheidet, ob es zur Ausf�hrung des aktuellen Branches kommt \cite{korel1990dynamic}. Sofern die
aktuelle Eingabe nicht zur gew�nschten Node f�hrt, wird ein neuer Eingabewert generiert \cite{korel1990dynamic}.

//Vergleich zu \cite{korel1990automated}, was ist anders?

Damit Informationen �ber gewisse Branches zur Verf�gung stehen, m�ssen diese klassifiziert werden. Dazu unterteilt Korel
diese in 4 Kategorien:

\begin{enumerate}
    \item A critical branch.
    \item A required branch.
    \item A semi-critical branch.
    \item A non-essential branch.
\end{enumerate}

Somit braucht man die \textit{critical branches} nicht verwerten, da diese nicht zur gew�nschten Node Y f�hren. Interessanter
sind die \textit{required branches}, da diese n�her an die gew�nschte Node f�hren k�nnen \cite{korel1990dynamic}. Diese werden
also nicht verworfen. Au�erdem gibt es noch \textit{semi-critical branches}, welche Einfluss auf den Weg zur Node Y haben
und in Zyklen enden k�nnen. Diese werden f�r den Suchalgorithmus verworfen. Die \textit{non-essential branches} werden
nicht verworfen, da diese keinen Einfluss auf die Node Y haben k�nnen \cite{korel1990dynamic}.

\section{Solving Subgoals}
Mit dem L�sen von sogenannten Subgoals werden auf dem Pfad zur gew�nschten Node jegliche Pfadbeschr�nkungen gel�st.
Auf dem Weg von einer Start-Node zu einer Node Y, durchl�uft das Programm einen entsprechenden Pfad. Sollte das Programm
nun einen Weg einschlagen, wodurch es die Ziel-Node Y nicht mehr erreichen kann, wird versucht ein neuen Eingabewert
zu generieren. Dieser neue Eingabewert soll im Idealfall an dieser Beschr�nkung vorbeikommen, sodass es zur gew�nschten
Node Y gelangt \cite{korel1990dynamic}. Korel beschreibt dies als \textit{Solving Subgoals}. Hierbei wird versucht mithilfe
von neuen zuf�lligen Eingabewerten die Beschr�nkung aufzul�sen.

Um dies zu vereinfachen, geht Korel davon aus, dass es sich hierbei um einfache Gleichungen beziehungsweise Ungleichungen
handelt \cite{korel1990dynamic}. Somit ist \(E_1\) \fett{op} \(E_2\) die Form der Ausdr�cke, wobei \(E_1\) und \(E_2\)
simple arithmetische Ausdr�cke sind und \fett{op} ein Vergleichsoperator ist \cite{korel1990dynamic}. Diese Ausdr�cke
k�nnen nach Tabelle \ref{tab:bspTabelle} in folgedene Formen umgewandelt werden: \(F\) \fett{rel} 0.

\begin{table}[h]
    \centering
    \caption{Predicate zu Function Relation nach Korel \cite{korel1990dynamic}}
    \label{tab:bspTabelle}
    \begin{tabular}{|c|c|c|} \hline
    Branch Predicate & Branch Function \(F\) & \fett{rel} \\ \hline
    $E_1 > E_2$ & $E_2 - E_1$ & $<$ \\ \hline
    $E_1 \geq E_2$ & $E_2 - E_1$ & $\leq$ \\ \hline
    $E_1 < E_2$ & $E_1 - E_2$ & $<$ \\ \hline
    $E_1 \leq E_2$ & $E_1 - E_2$ & $\leq$ \\ \hline
    $E_1 = E_2$ & abs($E_1 - E_2$) & $=$ \\ \hline
    $E_1 \neq E_2$ & abs($E_1 - E_2$) & $\leq$ \\ \hline
    \end{tabular}
\end{table}

Stelle man sich nun einen (zuf�lligen) Eingabewert \(x_0\) vor, mit dem das zu testende Programm gestartet wird, kann es
dazu kommen, dass die gew�nschte Node Y entweder erreicht wird oder nicht. Sollte die gew�nschte Node nicht erreicht werden
k�nnen, muss das erste \textit{Subgoal} gel�st werden \cite{korel1990dynamic}. Korel sagt nun, dass man sich eine Branch
Funktion \(F_j(x)\) vorstellen kann, welche genau die Pfadbeschr�nkung darstellt, welche das Programm nicht zur gew�nschten
Node gelangen l�sst \cite{korel1990dynamic}. Somit wird ein Wert \(x \in D\) gesucht, welcher $F_j(x)$  $\fett{rel}_j$  $0$
befriedigt, sodass die Pfadbeschr�nkung an dieser Stelle gel�st ist. Abh�ngig von $\fett{rel}_j$ kann man also eine
Funktion $F_j(x)$ minimieren und solange nach einer L�sung suchen bis $F_j(x)$ negativ wird \cite{korel1990dynamic}.
Die einfachste Art zu �berpr�fen, ob die Berechnung erfolgreich war, ist es das Programm mit der neuen Eingabe einfach
auszuf�hren. Wenn der gew�nschte Pfad nicht abgegangen wird, bezeichnet Korel dies als Verletzung des \textit{Constraint}
\cite{korel1990dynamic}.

\section{Probleme}

......................................................................
DART attempts to cover all executable program paths \cite{godefroid2005dart}

Dynamic Approach: geht nur in eine Richtung, Fokus liegt auf das Finden zu Node Y

dynamic test generation does not deal with functions calls, unknown code segments (such as library functions),

how to check at run-time whether predictions about new test inputs are matched in the next run, and does
not discuss completeness. \cite{godefroid2005dart}
.......................................................................
\section{Test Data Generation for Distributed Software}

\chapter{AUSTIN}

\chapter{EvoSuite}

\chapter{Schlussbemerkungen}
\section{Fazit}
\section{Ausblick}


\end{spacing}

% Die Inhalte des Anhangs werden analog zu den Kapiteln eingebunden.
% Wenn kein Anhang vorhanden ist, m�ssen die n�chten 7 Zeilen auskommentiert
% werden (bis '\end{spaching}').
\begin{spacing}{\zeilenabstandAnhang}
\appendix
    \chapter{Anhang}
    \label{sec:Anhang}
    % Hier k�nnen die einzelnen Anh�nge inkludiert werden. 
% Die Dateien m�ssen auf .tex enden. Diese Endung muss
% beim inkludieren aber weggelassen werden.
% Info: \include und \input unterscheiden sich im wesentlichen darin, dass bei
% \include immer eine neue Seite angefangen wird.
\section{Beispiele}
Zitat ohne Seitenangabe: \cite{BoeckleKnauberPohl+04}

Zitat mit Seitenangabe: \cite[S.1]{BoeckleKnauberPohl+04}

Referenz eines Glossareintrags: \gls{computer}

Eine normale Liste:
\begin{itemize}
\item Ein Punkt
\item Ein anderer Punkt
\end{itemize}

Eine nummerierte Liste:
\begin{enumerate}
\item Erstens
\item Zweitens
\item \ldots
\end{enumerate}

Eine Tabelle:
\begin{table}[h]
    \centering
    \captionabove{Eine Tabelle} 
    \label{tab:bspTabelle}
    \begin{tabular}{|c|c|c|} \hline
        Spalte 1 & Spalte 2 & Spalte 3 \\ \hline
        1 & 2 & 3 \\ \hline
    \end{tabular}
\end{table}

\textbf{Fetter Text}

\textit{Kursiver Text}

Eine Referenz: Siehe Tabelle \ref{tab:bspTabelle}: \nameref{tab:bspTabelle}
auf Seite \pageref{tab:bspTabelle}.

Ein Bild:
\begin{figure}[htb]
\centering
\includegraphics[width=0.2\textwidth]{Bilder/St_Uni-Logo-9-2003-eps-converted-to.pdf}
\caption{Das Logo der SUH}
\label{fig:suh-logo}
\end{figure}

Eine abgesetzte Formel:
\begin{equation}
\sum_{n=0}^{3}n=6
\end{equation}

Eine Formel im Flie�text: $2^2=4$ 

Ein Listing aus einer Datei:
\lstset{language=Java, basicstyle=\footnotesize, showstringspaces=false,tabsize=2}
\lstinputlisting[label=lst:HelloWorld,caption={HelloWorld}]{Listings/HelloWorld.java}

Ein 'on-the-fly' erstelltes Listing:
\begin{lstlisting}[label=lst:log,caption={Beispiel eines Log-Eintrags}] 
(Sun Sep 13 23:02:20 2009): ODBC Driver <system32>\wbemdr32.dll not present 
(Sun Sep 13 23:02:20 2009): Successfully verified WBEM OBDC adapter (incompatible version removed if it was detected).
(Sun Sep 13 23:02:20 2009): Wbemupgd.dll Registration completed.
\end{lstlisting}


 % Vor finaler Abgabe entfernen!

\end{spacing}

% Mit diesem Befehl werden nach nicht referenierte Glossareintr�ge
% ausgegeben (Nur zum Testen einkommentieren!)
% \glsaddall

% \addcontentsline{toc}{chapter}{Glossar}

% \ifthenelse{\boolean{final}}{\cleardoublepage}{\clearpage}
\clearpage
% \addcontentsline{toc}{chapter}{Bibliography}


% Literaturverzeichnis
%   Quelldatei ist: "Bibliographie.bib".
\bibliography{Bibliographie} % Aufruf: bibtex

% Verschiedene Presets f�r den Stil der Zitate und des Literaturverzeichnisses
% \bibliographystyle{plaindin} % Nur Zahlen, deutsch
% \bibliographystyle{plain} % Nur Zahlen, englisch
% \bibliographystyle{alphadin} % Alphanumerische K�rzel, deutsch
\bibliographystyle{alpha} % Alphanumerische K�rzel, englisch

\end{document}
